\begin{statement}{P3F.24}{Proved}
The forward sphere mapping torus is $(S^1\times S^6)\#\Sigma_\Phi$; both spin structures have Eells--Kuiper invariant $\mu(\Sigma_\Phi)$. A nonzero value excludes even an unoriented diffeomorphism with the product.
\end{statement}
\begin{proof}
Let \(\Phi\) be an orientation-preserving diffeomorphism of \(S^6\). Write \(\Sigma_\Phi=D^7\cup_\Phi(-D^7)\), with the direction of the boundary identification fixed from the first disc to the second. Its forward mapping torus satisfies \[
 T^+_\Phi=(S^6\times[0,1])/((z,1)\sim(\Phi(z),0))
       \cong(S^1\times S^6)\mathbin{\#}\Sigma_\Phi.   \tag{P3F.24.1}
\] Its reverse mapping torus, the convention used for the Jin six-dimensional mapping torus, satisfies \[
 T^-_\Phi=(S^6\times[0,1])/((z,0)\sim(\Phi(z),1))
       \cong(S^1\times S^6)\mathbin{\#}(-\Sigma_\Phi).\tag{P3F.24.2}
\] These are orientation-preserving diffeomorphisms for the indicated product and gluing orientations. For both spin structures on each mapping torus, its Eells--Kuiper invariant is respectively \(+\mu(\Sigma_\Phi)\) and \(-\mu(\Sigma_\Phi)\). In particular a nonzero such invariant excludes even an unoriented diffeomorphism with \(S^1\times S^6\).

\textbf{Proof.} First isotope \(\Phi\) to be the identity on a small oriented disc. The needed isotopy of disc embeddings is explicit: shrink the parametrized image disc toward its centre, move its centre along a path, move its oriented frame along the oriented frame bundle, and expand to the chosen target disc. Compactness gives a common sufficiently small radius along the path. Isotopy extension carries this isotopy to the ambient sphere. The result is a diffeomorphism equal to the identity on the target disc, and so supported in the complementary disc \(D^6\). Mapping tori of isotopic maps are diffeomorphic, by applying the isotopy along the interval; the same is true for their boundary-clutched spheres by extension over a collar. Thus this modification loses no data.

In the standard product \(S^1\times S^6\), choose a seam of the circle and the seven-ball \(B=D^6\times[-\epsilon,\epsilon]\) crossing that seam; round its corners. Its upper face has the positive orientation of \(D^6\), because the latter has even dimension. Define a boundary diffeomorphism of \(B\) to equal \(\Phi|_{D^6}\) on this face and the identity on the other face and the cylindrical collar. It is smooth since \(\Phi\) is the identity near \(\partial D^6\). Regluing this ball into its complement changes precisely the upper-to-lower seam identification to \(z\mapsto\Phi(z)\), and hence produces \(T^+_\Phi\). The boundary diffeomorphism, as an oriented map of \(\partial B\), is the original disc-supported \(\Phi\): both are its extension by the identity to the complementary six-disc.

The same regluing is connected sum with \(\Sigma_\Phi\). To check the map, remove an interior seven-ball from the second disc in \(D^7\cup_\Phi(-D^7)\). That remaining second piece is a collar \(S^6\times[0,1]\). Inserting this collar in the ball removed from \(S^1\times S^6\) gives exactly the specified boundary identification \(\Phi\). This proves (P3F.24.1) with its gluing direction. Replacing \(\Phi\) by \(\Phi^{-1}\) gives the reverse mapping torus and the reversed clutched sphere, proving \(P3F.24.2\).

There are exactly two spin structures because \(H^1(S^1\times S^6;\mathbb Z/2)=\mathbb Z/2\), and connected sum with a homotopy sphere does not change this group. Both product spin structures extend over \(S^1\times D^7\). That explicit eight-manifold has zero signature and zero first Pontryagin square. A homotopy seven-sphere bounds a spin eight-manifold; for instance the usual parallelizable representative supplied by \(\Theta _7=bP_8\) suffices. The boundary connected sum of this bounding manifold with \(S^1\times D^7\) bounds (P3F.24.1) for either spin structure. Its middle intersection form and its Pontryagin square are direct sums, since the joining handle has index one. Hence \[
 \mu(T^+_\Phi)=\mu(\Sigma_\Phi)
\] for each spin structure. The rational relative-to-absolute degree-four map is an isomorphism here because the mapping torus has \(H^3(-;\mathbb Q)=H^4(-;\mathbb Q)=0\). Thus the usual spin Eells--Kuiper expression \((p_1^2-4\sigma)/896\) applies without a choice of rational lift. Its invariance follows by gluing two such bounds: on the resulting closed spin eight-manifold this expression is the integral \(\widehat A\)-genus. Orientation reversal changes its sign. Both values for the product are zero, so neither preserving nor reversing orientation can identify a nonzero example with that product.
\end{proof}